% FAC-EN-U0065
% Source: fac_body.tex lines 2465--2483.
% Authority: PDF physical pp. 65--66; printed pp. 260--261
% (Chapter III, no. 67 complete).

\phantomsection\label{fac:III:3:no:67}
\textbf{67. A supplement.}

\phantomsection\label{fac:III:3:no:67:def:M-natural}
If $M$ is a graded $S$-module
\hyperref[fac:III:2:no:56:def:TF]{satisfying (TF)}, we shall denote by
$M^\natural$ the graded $S$-module $\Gamma(\mathcal{A}(M))$. We saw in
\phantomsection\label{fac:III:3:no:67:map:alpha}
\hyperref[fac:III:3:no:65:prop:5]{no.~65} that
$\alpha:M\longrightarrow M^\natural$ is
\hyperref[fac:III:2:no:56:def:C-morphism]{$\mathcal{C}$-bijective}. We shall
give conditions under which $\alpha$ is bijective.

\phantomsection\label{fac:III:3:no:67:prop:9}
\textsc{Proposition 9.} \textit{The homomorphism
$\alpha:M\longrightarrow M^\natural$ is bijective if and only if the
following conditions hold:}

\phantomsection\label{fac:III:3:no:67:prop:9:item:i}
\textit{(i) If $m\in M$ is such that $t_i\cdot m=0$ for every $i\in I$,
then $m=0$.}

\phantomsection\label{fac:III:3:no:67:prop:9:item:ii}
\textit{(ii) If homogeneous elements $m_i\in M$ of the same degree satisfy
$t_j\cdot m_i-t_i\cdot m_j=0$ for every pair $(i,j)$, then there exists
$m\in M$ such that $m_i=t_i\cdot m$.}

\phantomsection\label{fac:III:3:no:67:proof:necessity}
Let us show that conditions~\hyperref[fac:III:3:no:67:prop:9:item:i]{(i)}
and~\hyperref[fac:III:3:no:67:prop:9:item:ii]{(ii)} hold for $M^\natural$,
which will prove their necessity.
For~\hyperref[fac:III:3:no:67:prop:9:item:i]{(i)}, we may suppose that $m$ is
homogeneous, that is, that it is a section of $\mathcal{A}(M(n))$; in this
case, \hyperref[fac:III:3:no:67:prop:9:item:i]{the condition}
$t_i\cdot m=0$ implies that $m$ is zero on $U_i$, and,
since this holds for every $i\in I$, we indeed have $m=0$.
For~\hyperref[fac:III:3:no:67:prop:9:item:ii]{(ii)}, let
$n$ be the degree of the $m_i$; thus
$m_i\in\Gamma(\mathcal{A}(M(n)))$. Since $1/t_i$ is a section of
$\mathcal{O}(-1)$ on $U_i$, $m_i/t_i$ is a section of
$\mathcal{A}(M(n-1))$ on $U_i$, and
\hyperref[fac:III:3:no:67:prop:9:item:ii]{the condition}
$t_j\cdot m_i-t_i\cdot m_j=0$ shows that these various sections are the
\phantomsection\label{fac:III:3:no:67:proof:necessity:section-m}
restrictions of a unique section $m$ of $\mathcal{A}(M(n-1))$ on $X$. It
remains to compare the sections $t_i\cdot m$ and $m_i$; to show that they
coincide on $U_j$, it suffices to see that
$t_j(t_i\cdot m-m_i)=0$ on $U_j$, which follows from the formula
$t_j\cdot m_i=t_i\cdot m_j$ and
\hyperref[fac:III:3:no:67:proof:necessity:section-m]{the definition of $m$}.

\phantomsection\label{fac:III:3:no:67:proof:injectivity}
Let us now show that~\hyperref[fac:III:3:no:67:prop:9:item:i]{(i)} implies
that $\alpha$ is injective. For all sufficiently large $n$, we know that
$\alpha:M_n\longrightarrow M_n^\natural$ is bijective, and we may therefore
argue by descending induction on $n$. If $\alpha(m)=0$, with $m\in M_n$,
then
$t_i\alpha(m)=\alpha(t_i\cdot m)=0$. The induction hypothesis applies
because $t_i\cdot m\in M_{n+1}$, and gives $t_i\cdot m=0$ for every $i$;
condition~\hyperref[fac:III:3:no:67:prop:9:item:i]{(i)} then gives $m=0$.

\phantomsection\label{fac:III:3:no:67:proof:surjectivity}
Finally, let us show that~\hyperref[fac:III:3:no:67:prop:9:item:i]{(i)}
and~\hyperref[fac:III:3:no:67:prop:9:item:ii]{(ii)} imply that $\alpha$ is
surjective. As before, we may argue by descending induction on $n$. If
$m'\in M_n^\natural$, the induction hypothesis shows that there exist
$m_i\in M_{n+1}$ such that $\alpha(m_i)=t_i\cdot m'$. We have
$\alpha(t_j\cdot m_i-t_i\cdot m_j)=0$, whence
$t_j\cdot m_i-t_i\cdot m_j=0$, since $\alpha$ is injective.
Condition~\hyperref[fac:III:3:no:67:prop:9:item:ii]{(ii)} then implies the
existence of $m\in M_n$ such that $t_i\cdot m=m_i$; we have
$t_i(m'-\alpha(m))=0$ for every $i$, and
condition~\hyperref[fac:III:3:no:67:prop:9:item:i]{(i)}, already proved for
$M^\natural$, shows that $m'=\alpha(m)$ and completes the proof.

\phantomsection\label{fac:III:3:no:67:remark:1}
\textsc{Remarks.} (1) The proof shows that
condition~\hyperref[fac:III:3:no:67:prop:9:item:i]{(i)} is necessary and
sufficient for $\alpha$ to be injective.

\phantomsection\label{fac:III:3:no:67:remark:2}
(2) Conditions~\hyperref[fac:III:3:no:67:prop:9:item:i]{(i)}
and~\hyperref[fac:III:3:no:67:prop:9:item:ii]{(ii)} may be expressed as
follows:\hfill\break
\phantomsection\label{fac:III:3:no:67:formula:alpha1}the homomorphism
\mbox{$\alpha^1:M_n\longrightarrow H^0_1(M(n))$} is bijective for every
$n\in\mathbf{Z}$. Moreover,
\hyperref[fac:III:3:no:64:prop:4]{Proposition~4} shows that $M^\natural$ may
be identified with the graded $S$-module
\phantomsection\label{fac:III:3:no:67:formula:direct-sum}
$\sum_{n\in\mathbf{Z}}H^0(M(n))$, and it would be easy to deduce from this
a purely algebraic proof of
\hyperref[fac:III:3:no:67:prop:9]{Proposition~9} (without using the sheaf
$\mathcal{A}(M)$).
